%%%%%%%%%%%%%%%%%%%%%%%%%%%%%%%%%%%%%%%%%%%%%%%%%%%%%%%%%%%%%%%%%%%%%%
% Lớp 12 | Nguyên hàm, tích phân | Nguyên hàm
% Dạng 2D4?1-4: Nguyên hàm cơ bản hàm mũ, luỹ thừa
% Số câu: 62  (sinh tự động bởi tach_cau.py; từ nay thầy sửa trực tiếp file này)
%%%%%%%%%%%%%%%%%%%%%%%%%%%%%%%%%%%%%%%%%%%%%%%%%%%%%%%%%%%%%%%%%%%%%%

%%%%%%%%%%%%%%%% TRẮC NGHIỆM 4 PHƯƠNG ÁN %%%%%%%%%%%%%%%%

% nguồn: Phan1_Lop_12_Dai_so_Chuong_4
\begin{ex}%[2D4N1-4]%[ma:4beeaac5]
	Khẳng định nào sau đây  đúng?
	 \choice
		{\True  $\displaystyle\int 3^{2x}\mathrm{\,d}x=\dfrac{9^x}{\ln9}+C$}
		{$\displaystyle\int  3^{2x}\mathrm{\,d}x=9^x\cdot\ln9+C$}
		{$\displaystyle\int  3^{2x}\mathrm{\,d}x=\left(\dfrac{3^x}{\ln3}\right)^2+C$}
		 {$\displaystyle\int 3^{2x}\mathrm{\,d}x=3^x\cdot\ln3+C$}
	\loigiai{ 
		Ta có $\displaystyle\int 3^{2x}\mathrm{\,d}x=\displaystyle\int 9^{x}\mathrm{\,d}x=\dfrac{9^x}{\ln9}+C$ 
	 }
\end{ex}

% nguồn: Phan1_Lop_12_Dai_so_Chuong_4
\begin{ex}%[2D4N1-4]%[ma:e1b109f4]
	Gọi $F(x)$ là một nguyên hàm của hàm số $f(x)=\mathrm{e}^x$ thỏa mãn $F(0)=5$. Khẳng định nào sau đây là  đúng?
	\choice 
		 {$F(x)=\mathrm{e}^x$}
		 {\True $F(x)=\mathrm{e}^x+4$}
		{$F(x)=\mathrm{e}^x+5$} 
		{$F(x)=\mathrm{e}^x-4$} 
	\loigiai{Ta có 
		\[F(x)=\displaystyle \int \mathrm{e}^x\,\mathrm{\,d}x=  \mathrm{e}^x+C.\]
		Vì $F(0)=5$  nên
		\[F(0)=\mathrm{e}^0+C=5 \Leftrightarrow C=4.\] 
	Vậy ta  có $F(x)=\mathrm{e}^x+4$.}
\end{ex}

% nguồn: Phan1_Lop_12_Dai_so_Chuong_4
\begin{ex}%[2D4N1-4]%[ma:293b52ae]
	 Trong các khẳng định sau, khẳng định nào \textbf{sai}?
	 \choice
		 {$\displaystyle\int \mathrm{\,d} x=x+C$}
		{$\displaystyle\int \cos x  \mathrm{\,d} x=\sin x+C$}
		{\True $\displaystyle\int \mathrm{e}^{x} \mathrm{\,d} x=\dfrac{\mathrm{e}^{x+1}}{x+1}+C$} 
		 {$\displaystyle\int \dfrac{1}{x} \mathrm{\,d} x=\ln |x|+C$}
	 \loigiai{
		Ta có $\displaystyle\int \mathrm{e}^{x}  \mathrm{\,d} x=\mathrm{e}^{x}+C$, do đó $\displaystyle\int \mathrm{e}^{x} \mathrm{\,d} x=\dfrac{\mathrm{e}^{x+1}}{x+1}+C$ là sai.
	} 
\end{ex}

% nguồn: Phan1_Lop_12_Dai_so_Chuong_4
\begin{ex}%[2D4N1-4]%[ma:7ac76f49]
	Họ tất cả các nguyên hàm của hàm  số $f(x)=2026x^{2025}+\mathrm{e}^x$ là
	 \choice
		{$2026  x^{2026}+\mathrm{e}^x+C$}
		 {$x^{2026}+\mathrm{e}^x+C$}
		{$2026  x+\mathrm{e}^x+C$}
		 {\True $x^{2026}+\mathrm{e}^x+C$}
	\loigiai{ 
		Ta  có $\displaystyle\int\left(2026x^{2025}+\mathrm{e}^x\right)\mathrm{\,d}x=\dfrac{2026 x^{2026}}{2026}+\mathrm{e}^x+C=x^{2026}+\mathrm{e}^x+C$.
	 }
\end{ex}

% nguồn: Phan1_Lop_12_Dai_so_Chuong_4
\begin{ex}%[2D4N1-4]%[ma:64eaf71b]
	Tất cả các nguyên hàm  của hàm số $f(x)=3^{-x}$ là
	 \choice
		{\True $-\dfrac{3^{-x}}{\ln 3}+C$} 
		 {$-{3^{-x}}+C$}
		 {${3^{-x}} \ln 3 +C$}
		{$\dfrac{3^{-x}}{\ln  3}+C$}
	\loigiai{Tất cả các nguyên hàm của hàm  số $f(x)=3^{-x}$ là $-\dfrac{3^{-x}}{\ln 3}+C$.}
\end{ex}

% nguồn: Phan1_Lop_12_Dai_so_Chuong_4
\begin{ex}%[2D4N1-4]%[ma:1311c02f]
	Cho hàm số $f(x)=2^x+3$. Khẳng định nào sau  đây đúng?
	 \choice
		{$\displaystyle\int f(x)\, \mathrm{d}x=2^{x-3}+C$} 
		{$\displaystyle\int f(x)\, \mathrm{d}x=2^x \ln  2+3x+C$}
		{$\displaystyle\int f(x)\, \mathrm{d}x=\dfrac{2^x}{\ln 2}+3+C$} 
		{\True $\displaystyle\int f(x)\, \mathrm{d}x=\dfrac{2^x}{\ln  2}+3x+C$}
	\loigiai{ 
		Ta  có
		$\displaystyle\int f(x)\, \mathrm{d}x=\displaystyle\int (2^x+3)\, \mathrm{d}x=\displaystyle\int 2^x\, \mathrm{d}x+\displaystyle\int 3\, \mathrm{d}x=\dfrac{2^x}{\ln 2}+3x+C$. 
	} 
\end{ex}

% nguồn: Phan1_Lop_12_Dai_so_Chuong_4
\begin{ex}%[2D4N1-4]%[ma:4a6e9301]
	Cho $a$  là số dương khác 1. Hàm số $y = a^x$ là một nguyên hàm của hàm số
	 \choice
		{$y  = a^x$}
		 {$y = a^{x+1}$}
		{$y =  \dfrac{a^x}{\ln a}$}
		{\True $y = a^x \ln a$} 
	\loigiai{ 
		Ta  có $y' = (a^x)'=a^x \ln a$.
	} 
\end{ex}

% nguồn: Phan1_Lop_12_Dai_so_Chuong_4
\begin{ex}%[2D4N1-4]%[ma:a85e6983]
	Một  nguyên hàm của hàm số $f(x)=2\,025^x$ là
	\choice 
		{$F(x)=x \cdot  2\,025^{x-1}$}
		{$F(x)=2\,025  \cdot 2\,024^x$}
		{$F(x)=2\,025^x \cdot \ln 2\,025$} 
		{\True $F(x)=\dfrac{2\,025^x}{\ln 2\,025}$} 
	\loigiai{ 
		 Ta có
		\[\displaystyle\int f(x)\mathrm{\,d}x=\displaystyle\int 2\,025^x\mathrm{\,d}x=\dfrac{2\,025^x}{\ln 2\,025}+C.\] 
	} 
\end{ex}

% nguồn: Phan1_Lop_12_Dai_so_Chuong_4
\begin{ex}%[2D4N1-4]%[ma:006847d2]
	Tìm nguyên hàm của hàm số $f(x)=2025^{x}$ 
	\choice 
		{$2025^{x} \ln 2025+C$} 
		{$2025^{x+1}+C$} 
		{$\dfrac{2025^{x+1}}{x+1}+C$} 
		{\True $\dfrac{2025^{x}}{\ln 2025}+C$} 
	\loigiai{ 
		Áp dụng công thức $\displaystyle\int a^{x} \mathrm{\,d}  x=\dfrac{a^{x}}{\ln a}+C$, $(0<a \neq 1)$.
	 }
\end{ex}

% nguồn: Phan1_Lop_12_Dai_so_Chuong_4
\begin{ex}%[2D4N1-4]%[ma:266d3b40]
	Tìm nguyên hàm của hàm số  $f(x)=3^x$.
	 \choice
		{\True $\displaystyle\int 3^x\mathrm{\,d}x=\dfrac{3^x}{\ln 3}+C$} 
		{$\displaystyle\int  3^x\mathrm{\,d}x=3^x\ln 3+C$}
		{$\displaystyle\int 3^x\mathrm{\,d}x=3^{x+1}+C$} 
		 {$\displaystyle\int 3^x\mathrm{\,d}x=\dfrac{3^{x+1}}{x+1}+C$}
	 \loigiai{
		Ta có $\displaystyle\int a^x\mathrm{\,d}x=\dfrac{a^x}{\ln  a}+C$ nên $\displaystyle\int 3^x\mathrm{\,d}x=\dfrac{3^x}{\ln 3}+C$.
	} 
\end{ex}

% nguồn: Phan1_Lop_12_Dai_so_Chuong_4
\begin{ex}%[2D4N1-4]%[ma:0adf1a93]
	Họ nguyên hàm của hàm  số $f(x)=7^x$ là
	 \choice
		 {$7^x+C$}
		{$7^x  \ln 7+C$}
		 {\True $\dfrac{7^x}{\ln 7}+C$}
		 {$\dfrac{7^{x+1}}{x+1}+C$}
	\loigiai{ 
		Ta có $\displaystyle\int 7^x\mathrm{\,d}x=\dfrac{7^x}{\ln7}+C$. 
	 }
\end{ex}

% nguồn: Phan1_Lop_12_Dai_so_Chuong_4
\begin{ex}%[2D4N1-4]%[ma:3d0c048e]
	Khẳng  định nào dưới đây là sai?
	 \choice
		{$\displaystyle\int x^n\mathrm{\,d}x=\dfrac{x^{n+1}}{n+1}+C$,  $n\neq -1$}
		 {$\displaystyle\int \mathrm{e}^x\mathrm{\,d}x=\mathrm{e}^x+C$}
		{$\displaystyle\int \cos x\mathrm{\,d}x=\sin x+C$} 
		{\True  $\displaystyle\int 17^x\mathrm{\,d}x=\dfrac{17^{x+1}}{x+1}+C$}
	\loigiai{ 
		 Ta có $\displaystyle\int 17^x\mathrm{\,d}x=\dfrac{17^x}{\ln 17}+C$.
	 }
\end{ex}

% nguồn: Phan1_Lop_12_Dai_so_Chuong_4
\begin{ex}%[2D4N1-4]%[ma:1f738c11]
	Họ  nguyên hàm của hàm số $f(x)=\mathrm{e}^{x}+2 x$ là
	 \choice
		{$\dfrac{1}{x+1} \mathrm{e}^{x}+x^{2}+C$} 
		 {$\mathrm{e}^{x}-x^{2}+C$}
		{$\mathrm{e}^{x}-2+C$} 
		 {\True $\mathrm{e}^{x}+x^{2}+C$}
	 \loigiai{}
\end{ex}

% nguồn: Phan1_Lop_12_Dai_so_Chuong_4
\begin{ex}%[2D4N1-4]%[ma:ac5bb642]
	Tìm $\displaystyle\int\mathrm{e}^x\mathrm{\,d}x$. 
	\choice 
		 {$\displaystyle\int\mathrm{e}^x\mathrm{\,d}x=\dfrac{1}{x\ln 2}+C$}
		{\True  $\displaystyle\int\mathrm{e}^x\mathrm{\,d}x=\mathrm{e}^x+C$}
		 {$\displaystyle\int\mathrm{e}^x\mathrm{\,d}x=\mathrm{e}^x\ln 2+C$}
		{$\displaystyle\int\mathrm{e}^x\mathrm{\,d}x=\dfrac{1}{\mathrm{e}^x}+C$} 
	 \loigiai{Ta có $\displaystyle\int\mathrm{e}^x\mathrm{\,d}x=\mathrm{e}^x+C$.}
\end{ex}

% nguồn: Phan1_Lop_12_Dai_so_Chuong_4
\begin{ex}%[2D4N1-4]%[ma:f9b9637b]
	Nguyên hàm của hàm số $f(x)=2x+6^x$  là
	 \choice
		 {$\displaystyle\int f(x)\mathrm{\,d}x=2+6^x\ln 6+C$}
		{$\displaystyle\int f(x)\mathrm{\,d}x=x^2+\dfrac{6^{x+1}}{x+1}+C$} 
		{\True $\displaystyle\int f(x)\mathrm{\,d}x=x^2+\dfrac{6^x}{\ln  6}+C$}
		{$\displaystyle\int  f(x)\mathrm{\,d}x=2x^2+6^x\ln 6+C$}
	 \loigiai{
		Ta có $\displaystyle\int f(x)\mathrm{\,d}x=x^2+\dfrac{6^x}{\ln  6}+C$.
	 }
\end{ex}

% nguồn: Phan1_Lop_12_Dai_so_Chuong_4
\begin{ex}%[2D4N1-4]%[ma:69f2a72d]
	Cho $a$, $b \in \mathbb{R}$, $a \ne 0$. Hàm số $y = \mathrm{e}^{ax+b}$ là một  nguyên hàm của hàm số
	\choice 
		{$y  = -\mathrm{e}^{ax+b}$}
		{$y =  \mathrm{e}^{ax+b}$}
		{\True  $y = a \mathrm{e}^{ax+b}$}
		{$y  = \dfrac{\mathrm{e}^{ax+b}}{a}$}
	 \loigiai{
		Ta có $y'  = (\mathrm{e}^{ax+b})'=a\cdot \mathrm{e}^{ax+b}$.
	 }
\end{ex}

% nguồn: Phan1_Lop_12_Dai_so_Chuong_4
\begin{ex}%[2D4N1-4]%[ma:0e0d0b34]
	Nguyên hàm của  hàm số $f(x)=8^{x}$ là
	 \choice
		 {$\dfrac{8^{x+1}}{x+1}+C$}
		{\True $\dfrac{8^{x}}{\ln  8}+C$}
		{$x8^{x-1}+C$} 
		{$8^{x}\ln  8+C$}
	\loigiai{ 
		Nguyên hàm của  hàm số $f(x)=8^x$ là $\dfrac{8^{x}}{\ln 8}+C$.
	 }
\end{ex}

% nguồn: Phan1_Lop_12_Dai_so_Chuong_4
\begin{ex}%[2D4H1-4]%[ma:b4deeca0]
	Phát biểu  nào sau đây là đúng?
	\choice 
		{$\displaystyle\int  12^x\mathrm{\,d}x=12^x$}
		{$\displaystyle\int 12^x\mathrm{\,d}x=\dfrac{12^x}{\ln  12}$}
		{\True $\displaystyle\int 12^x\mathrm{\,d}x=\dfrac{12^x}{\ln 12}+C$} 
		{$\displaystyle\int 12^x\mathrm{\,d}x=12^x+C$} 
	\loigiai{ 
		Ta có $\displaystyle\int 12^x\mathrm{\,d}x=\dfrac{12^x}{\ln 12}+C$. 
	} 
\end{ex}

% nguồn: Phan1_Lop_12_Dai_so_Chuong_4
\begin{ex}%[2D4H1-4]%[ma:681b1faa]
	Cho hàm số $f(x)$ thoả mãn $f'(x)=2x-3e^x, \forall x \in \mathbb{R}$ và $f(0)=5$. Khẳng định  nào sau đây \textbf{đúng}?
	\choice 
		 {$f(x)=2x+3e^x+2$}
		 {$f(x)=x^2+3e^x+2$}
		{$f(x)=2x-3e^x+8$} 
		 {\True $f(x)=x^2-3e^x+8$}
	\loigiai{ 
		Ta có $f(x)=\displaystyle\int  \left(2x-3e^x\right) \mathrm{\,d}x=x^2-3e^x+C$.\\
		Có $f(0)=5\Rightarrow C=8$.  Vậy $f(x)=x^2-3e^x+8$.
	 }
\end{ex}

% nguồn: Phan1_Lop_12_Dai_so_Chuong_4
\begin{ex}%[2D4H1-4]%[ma:cbc18dac]
	Cho $y=f(x)$ có đạo  hàm $f'(x)=\mathrm{e}^x+2x+1$ và $f(0)=1$. Biết $F(x)$ là nguyên hàm của $f(x)$ thỏa mãn $F(1)=e$. Tính $F(0)$.
	 \choice
		 {$\dfrac{5}{6}$}
		 {$-\dfrac{1}{6}$}
		 {\True $\dfrac{1}{6}$}
		{$-\dfrac{5}{6}$} 
	 \loigiai{
		Ta có \[f'(x)=\mathrm{e}^x+2x+1 \Rightarrow f(x)=\displaystyle\int  f'(x)\mathrm{\,d}x=\mathrm{e}^x+x^2+x+C.\]
		 Do
		\[f(0)=1 \Rightarrow \mathrm{e}^0+0+0+C=1 \Rightarrow C=0 \Rightarrow  f(x)=\mathrm{e}^x+x^2+x.\]
		Suy  ra
		\[F(x)=\displaystyle\int f(x)\mathrm{\,d}x=\mathrm{e}^x+\dfrac{x^3}{3}+\dfrac{x^2}{2}+C.\] 
		Khi đó 
		\[F(1)=e \Rightarrow  e+\dfrac{1}{3}+\dfrac{1}{2}+C=e \Rightarrow C=-\dfrac{5}{6}.\]
		Vậy $F(0)=1+0+0-\dfrac{5}{6}=\dfrac{1}{6}$. 
	 }
\end{ex}

% nguồn: Phan1_Lop_12_Dai_so_Chuong_4
\begin{ex}%[2D4H1-4]%[ma:9d57b81d]
	Khẳng định nào  sau đây \textbf{đúng}?
	\choice 
		{$\displaystyle\int 3^{2x} \mathrm{\,d}x=\dfrac{3^{2x}}{\ln 3}+C$} 
		 {\True $\displaystyle\int 3^{2x} \mathrm{\,d}x=\dfrac{9^x}{\ln 9}+C$}
		{$\displaystyle\int 3^{2x}  \mathrm{\,d}x=3^{2x} \cdot \ln 3+C$}
		{$\displaystyle\int 3^{2x}  \mathrm{\,d}x=9^x \cdot \ln 9+C$}
	 \loigiai{
		Ta có $\displaystyle\int 3^{2x} \mathrm{\,d}x=\displaystyle\int 9^{x} \mathrm{\,d}x=\dfrac{9^x}{\ln  9}+C$.
	 }
\end{ex}

% nguồn: Phan1_Lop_12_Dai_so_Chuong_4
\begin{ex}%[2D4H1-4]%[ma:c36ecb8b]
	Hàm số nào sau đây  là một nguyên hàm của hàm số $f(x) = 2^x + \cos x$?
	\choice 
		 {$F(x) = 2^x \ln 2 + \sin x$}
		 {\True $F(x) = \dfrac{2^x}{\ln 2} + \sin x$}
		{$F(x) = \dfrac{2^x}{\ln 2} - \sin x$} 
		{$F(x) = 2^x \ln  2 - \sin x$}
	\loigiai{ 
	Ta  có:
	 \[
	\int (2^x + \cos x) \mathrm{\,d}x =  \dfrac{2^x}{\ln 2} + \sin x + C.
	 \]
	So sánh  với các đáp án, hàm số $F(x) = \dfrac{2^x}{\ln 2} + \sin x$ là một nguyên hàm của $f(x)$.
	} 
\end{ex}

% nguồn: Phan1_Lop_12_Dai_so_Chuong_4
\begin{ex}%[2D4H1-4]%[ma:9586f489]
	Họ nguyên hàm của hàm  số $f(x)=4^x+\cos 2x$ là
	\choice 
		{$\dfrac{4^x}{\ln 4}-\dfrac{\sin 2x}{2}+C$} 
		{$4^x  \ln 4+\dfrac{\sin 2x}{2}+C$}
		{$4^x  \ln 4-\dfrac{\sin 2x}{2}+C$}
		{\True $\dfrac{4^x}{\ln 4}+\dfrac{\sin 2x}{2}+C$} 
	 \loigiai{
		Ta có  \[\displaystyle\int f(x)\mathrm{\,d}x=\displaystyle\int 4^x\mathrm{\,d}x+\displaystyle\int \cos 2x\mathrm{\,d}x=\dfrac{4^x}{\ln 4}+\dfrac{\sin 2x}{2}+C.\]
	} 
\end{ex}

% nguồn: Phan1_Lop_12_Dai_so_Chuong_4
\begin{ex}%[2D4H1-4]%[ma:033d2e4a]
	Một nguyên hàm của  hàm số $f(x)=\mathrm{e}^{3x+1}-2x^2$ là
	 \choice
		 {$\dfrac{\mathrm{e}^{3x+1}-2x^3}{3}$}
		 {$\dfrac{\mathrm{e}^{3x+1}}{3}-x^3$}
		 {\True $\dfrac{\mathrm{e}^{3x+1}}{3}-2x^3$}
		 {$\dfrac{\mathrm{e}^{3x+1}-x^3}{3}$}
	 \loigiai{
		Ta có \[\displaystyle\int f(x)\mathrm{\,d}x=\displaystyle\int \mathrm{e}^{3x+1}\mathrm{\,d}x-\displaystyle\int  2x^2\mathrm{\,d}x=\dfrac{\mathrm{e}^{3x+1}}{3}-\dfrac{2x^3}{3}+C=\dfrac{\mathrm{e}^{3x+1}-2x^3}{3}+C.\]
	} 
\end{ex}

% nguồn: Phan1_Lop_12_Dai_so_Chuong_4
\begin{ex}%[2D4H1-4]%[ma:4bb6c62a]
	Tìm $F(x)$ là một nguyên hàm của hàm số $f(x)=\mathrm{e}^x-1$ trên $(-\infty;+\infty)$, biết  $F(0)=2$.
	 \choice
		 {$F(x)=\ln x-x-1$}
		 {$F(x)=\mathrm{e}^x-x-1$}
		{$F(x)=\dfrac{1}{\mathrm{e}^x}-x+1$} 
		{\True  $F(x)=\mathrm{e}^x-x+1$}
	 \loigiai{
		 Ta có \[\displaystyle\int (\mathrm{e}^x-1)\mathrm{\,d}x=\mathrm{e}^x-x+C.\]
		Do  $F(0)=2 \Rightarrow 1-0+C=2 \Rightarrow C=1$.\\
		Vậy  $F(x)=\mathrm{e}^x-x+1$.
	} 
\end{ex}

% nguồn: Phan1_Lop_12_Dai_so_Chuong_4
\begin{ex}%[2D4H1-4]%[ma:4a6638e1]
	 Nguyên hàm của hàm số $f(x)=2^x(1+2^{-x} \cdot \sin x)$ là
	\choice 
		{$\dfrac{2^{x+1}}{x+1}-\cos x+C$} 
		{\True $\dfrac{2^x}{\ln 2}-\cos x+C$} 
		 {$\dfrac{2^x}{\ln 2}+\cos x+C$}
		{$\dfrac{2^{x-1}}{x+1}+\cos  x+C$}
	 \loigiai{
		Ta có \[f(x)=2^x+\sin  x \Rightarrow \displaystyle\int f(x)\mathrm{\,d}x=\dfrac{2^x}{\ln 2}-\cos x+C.\]
	 }
\end{ex}

% nguồn: Phan1_Lop_12_Dai_so_Chuong_4
\begin{ex}%[2D4H1-4]%[ma:35cfed20]
	Tìm nguyên hàm  của hàm số $f(x)=\mathrm{e}^{-x}\left(2+\dfrac{\mathrm{e}^x}{\cos^2 x}\right)$
	 \choice
		{\True $F(x)=-\dfrac{2}{\mathrm{e}^x}+\tan x+C$} 
		{$F(x)=2\mathrm{e}^x-\tan x+C$} 
		{$F(x)=-\dfrac{2}{\mathrm{e}^x}-\tan x+C$} 
		 {$F(x)=2\mathrm{e}^{-x}+\tan x+C$}
	\loigiai{ 
		 Ta có \[f(x)=2\mathrm{e}^{-x}+\dfrac{1}{\cos^2 x}=2\mathrm{e}^{-x}+\tan' x.\]
		 Vậy \[\displaystyle\int f(x)\mathrm{\,d}x=-2\mathrm{e}^{-x}+\tan x+C=-\dfrac{2}{\mathrm{e}^x}+\tan x+C.\]
	 }
\end{ex}

% nguồn: Phan1_Lop_12_Dai_so_Chuong_4
\begin{ex}%[2D4H1-4]%[ma:4a0f04cf]
	Hàm số $F(x)=x^2+\mathrm{e}^{-2x}+2\,025$ là một nguyên hàm của hàm số nào  dưới đây?
	\choice 
		{$f(x)=\dfrac{x^3}{3}-\dfrac{1}{2}\mathrm{e}^{-2x}+2\,025x$} 
		{\True $f(x)=2x-2\mathrm{e}^{-2x}$} 
		 {$f(x)=2x-2\mathrm{e}^{-2x}+2\,025x$}
		{$f(x)=2x+2\mathrm{e}^{-2x}+2\,025$} 
	 \loigiai{
		 Ta có \[F'(x)=(x^2+\mathrm{e}^{-2x}+2\,025)'=2x-2\mathrm{e}^{-2x}.\]
	} 
\end{ex}

% nguồn: Phan1_Lop_12_Dai_so_Chuong_4
\begin{ex}%[2D4H1-4]%[ma:762deb74]
	Cho  hàm số $f(x)=1+\mathrm{e}^x$. Khẳng định nào dưới đây đúng?
	 \choice
		 {$\displaystyle\int f(x)\mathrm{\,d}x=x+x \mathrm{e}^{x-1}+C$}
		{$\displaystyle\int f(x)\mathrm{\,d}x=1+\mathrm{e}^x+C$} 
		{\True $\displaystyle\int f(x)\mathrm{\,d}x=x+\mathrm{e}^x+C$} 
		{$\displaystyle\int  f(x)\mathrm{\,d}x=1+x \mathrm{e}^{x-1}+C$}
	\loigiai{ 
		Ta có \[\displaystyle\int f(x)\mathrm{\,d}x=\displaystyle\int (1+\mathrm{e}^x)\mathrm{\,d}x=x+\mathrm{e}^x+C.\] 
	 }
\end{ex}

% nguồn: Phan1_Lop_12_Dai_so_Chuong_4
\begin{ex}%[2D4H1-4]%[ma:b4521ef3]
	 Hàm số $F(x)=\dfrac{x^3}{3}+\mathrm{e}^x$ là một nguyên hàm của hàm số nào sau đây?
	\choice 
		{$f(x)=\dfrac{x^4}{3}+\mathrm{e}^x$} 
		{$f(x)=3x^2+\mathrm{e}^x$} 
		 {$f(x)=\dfrac{x^4}{12}+\mathrm{e}^x$}
		{\True  $f(x)=x^2+\mathrm{e}^x$}
	 \loigiai{
		Ta có \[f(x)=F'(x)=x^2+\mathrm{e}^x.\] 
	} 
\end{ex}

% nguồn: Phan1_Lop_12_Dai_so_Chuong_4
\begin{ex}%[2D4H1-4]%[ma:3a7a77b2]
	Nguyên hàm $F(x)$ của  hàm số $f(x)=\mathrm{e}^x+2\mathrm{e}^{-x}-2x$ thỏa mãn $F(0)=1$
	\choice 
		 {\True $F(x)=\mathrm{e}^x-2\mathrm{e}^{-x}-x^2+2$}
		{$F(x)=\mathrm{e}^x+2\mathrm{e}^{-x}-x^2-2$} 
		{$F(x)=\mathrm{e}^x-2\mathrm{e}^{-x}-x^2$} 
		 {$F(x)=\mathrm{e}^x+2\mathrm{e}^{-x}-x^2$}
	\loigiai{ 
		 Ta có \[\displaystyle\int f(x)\mathrm{\,d}x=\displaystyle\int (\mathrm{e}^x+2\mathrm{e}^{-x}-2x)\mathrm{\,d}x=\mathrm{e}^x-2\mathrm{e}^{-x}-x^2+C.\]
		Vì $F(0)=1 \Rightarrow  1-2+C=1 \Rightarrow C=2$.\\
		Vậy $F(x)=\mathrm{e}^x-2\mathrm{e}^{-x}-x^2+2$. 
	 }
\end{ex}

% nguồn: Phan1_Lop_12_Dai_so_Chuong_4
\begin{ex}%[2D4H1-4]%[ma:d80e70b0]
	 Biết $F(x)$ là một nguyên hàm của $f(x)=\mathrm{e}^{2x}$ và $F(0)=0$. Giá trị $F(\ln 3)$ bằng
	 \choice
		{$2$} 
		 {$6$}
		 {$8$}
		{\True $4$} 
	\loigiai{ 
		Ta có\\ 
		$\displaystyle\int  \mathrm{e}^{2x}\mathrm{\,d}x=\dfrac{1}{2}\mathrm{e}^{2x}+C$.\\
		$F(0)=0 \Rightarrow \dfrac{1}{2}\mathrm{e}^0+C=0 \Rightarrow C=-\dfrac{1}{2}$.\\ 
		$F(\ln 3)=\dfrac{1}{2}\mathrm{e}^{2\ln 3}-\dfrac{1}{2}=\dfrac{1}{2} \cdot  9-\dfrac{1}{2}=4$.
	 }
\end{ex}

% nguồn: Phan1_Lop_12_Dai_so_Chuong_4
\begin{ex}%[2D4H1-4]%[ma:7d8dc057]
	Cho hàm số $f(x)=2^x+3$. Khẳng  định nào sau đây đúng?
	\choice 
		 {$\displaystyle\int f(x)\mathrm{\,d}x=2^{x-3}+C$}
		{$\displaystyle\int f(x)\mathrm{\,d}x=2^x \ln  2+3x+C$}
		 {$\displaystyle\int f(x)\mathrm{\,d}x=\dfrac{2^x}{\ln 2}+3+C$}
		{\True $\displaystyle\int  f(x)\mathrm{\,d}x=\dfrac{2^x}{\ln 2}+3x+C$}
	\loigiai{ 
		Ta có \[\displaystyle\int f(x)\mathrm{\,d}x=\displaystyle\int 2^x\mathrm{\,d}x+\displaystyle\int 3\mathrm{\,d}x=\dfrac{2^x}{\ln 2}+3x+C.\] 
	} 
\end{ex}

% nguồn: Phan1_Lop_12_Dai_so_Chuong_4
\begin{ex}%[2D4H1-4]%[ma:11f5c173]
	Biết hàm số $F(x)$ là một nguyên hàm của hàm số $f(x)=4^x$ và $F(1)=\dfrac{1}{\ln 2}$. Giá trị của  $F\left(\dfrac{3}{2}\right) \cdot \ln 2$ bằng
	\choice 
		 {\True $3$}
		 {$7$}
		{$5$} 
		 {$9$}
	\loigiai{ 
		Ta có $F(x)=\displaystyle\int 4^x \mathrm{\,d}x=\dfrac{4^x}{2\ln  2}+C$.\\
		Có $F(1)=\dfrac{1}{\ln 2}\Rightarrow  C=-\dfrac{1}{\ln 2}$. Suy ra $F(x)=\dfrac{4^x}{2\ln 2}-\dfrac{1}{\ln 2}$.\\
		Vậy  $F\left(\dfrac{3}{2}\right) \cdot \ln 2=3$.
	 }
\end{ex}

% nguồn: Phan1_Lop_12_Dai_so_Chuong_4
\begin{ex}%[2D4H1-4]%[ma:4ea1c0df]
	Trong các hàm số dưới đây, hàm số nào  có họ tất cả các nguyên hàm là $F(x)=\dfrac{a^x}{\ln a}+C$ với $a>0$, $a\ne 1$?
	 \choice
		 {\True $f(x)=a^x$}
		{$f(x)=\dfrac{1}{x}$} 
		{$f(x)=\ln  x$}
		{$f(x)=x^a$} 
	\loigiai{ 
		Ta có \[\displaystyle\int a^x\mathrm{\,d}x=\dfrac{a^x}{\ln  a}+C.\]
	 }
\end{ex}

% nguồn: Phan1_Lop_12_Dai_so_Chuong_4
\begin{ex}%[2D4H1-4]%[ma:552c6ab0]
	Tìm một nguyên hàm $F(x)$ của hàm số $f(x)=\mathrm{e}^x-1$ trên $(-\infty; +\infty)$,  biết $F(0)=2$.
	 \choice
		{$F(x)=\dfrac{1}{\mathrm{e}^x}-x+1$} 
		{$F(x)=\ln x-x-1$} 
		 {$F(x)=\mathrm{e}^x-x-1$}
		 {\True $F(x)=\mathrm{e}^x-x+1$}
	 \loigiai{
		Ta có  \[\displaystyle\int (\mathrm{e}^x-1)\mathrm{\,d}x=\mathrm{e}^x-x+C.\]
		Vì $F(0)=2 \Rightarrow  \mathrm{e}^0-0+C=2 \Rightarrow C=1$.\\
		Vậy $F(x)=\mathrm{e}^x-x+1$. 
	} 
\end{ex}

% nguồn: Phan1_Lop_12_Dai_so_Chuong_4
\begin{ex}%[2D4H1-4]%[ma:51d6a3ae]
	Kết quả  $\displaystyle\int (x+\mathrm{e}^{2\,020x})\mathrm{\,d}x$ bằng
	 \choice
		{$x^2+\dfrac{\mathrm{e}^{2\,020x}}{2\,020}+C$} 
		{$x^3+\dfrac{\mathrm{e}^{2\,020x}}{2\,020}+C$} 
		 {\True $\dfrac{x^2}{2}+\dfrac{\mathrm{e}^{2\,020x}}{2\,020}+C$}
		 {$x+\dfrac{\mathrm{e}^{2\,020x}}{2\,020}+C$}
	\loigiai{ 
		Ta  có \[\displaystyle\int (x+\mathrm{e}^{2\,020x})\mathrm{\,d}x=\dfrac{x^2}{2}+\dfrac{\mathrm{e}^{2\,020x}}{2\,020}+C.\]
		Vậy nguyên hàm cần  tìm là $\dfrac{x^2}{2}+\dfrac{\mathrm{e}^{2\,020x}}{2\,020}+C$.
	 }
\end{ex}

% nguồn: Phan1_Lop_12_Dai_so_Chuong_4
\begin{ex}%[2D4H1-4]%[ma:834989da]
	Họ tất cả các nguyên hàm của  hàm số $f(x)=5^{2x-1}$ là
	 \choice
		{$\dfrac{1}{2} \cdot  5^{2x-1} \cdot \ln 5+C$}
		{$2 \cdot 5^{2x-1}  \cdot \ln 5+C$}
		{$5^{2x-1} \cdot \ln 5+C$} 
		{\True $\dfrac{1}{2\ln 5} \cdot  5^{2x-1}+C$}
	 \loigiai{
		 Ta có \[\displaystyle\int 5^{2x-1} \mathrm{\,d}x=\dfrac{1}{2\ln 5} \cdot 5^{2x-1}+C.\]
	} 
\end{ex}

% nguồn: Phan1_Lop_12_Dai_so_Chuong_4
\begin{ex}%[2D4H1-4]%[ma:c95ec96b]
	Cho hàm số $f(x)=3^x+2x$. Khẳng định nào dưới đây đúng? 
	\choice 
		 {$\displaystyle\int f(x)\mathrm{\,d}x=\dfrac{3^x}{\ln 3}+2+C$}
		{\True $\displaystyle\int f(x)\mathrm{\,d}x=\dfrac{3^x}{\ln 3}+x^2+C$} 
		{$\displaystyle\int f(x)\mathrm{\,d}x=3^x  \cdot \ln 3+2+C$}
		 {$\displaystyle\int f(x)\mathrm{\,d}x=3^x \cdot \ln 3+x^2+C$}
	 \loigiai{
		Ta có \[\displaystyle\int f(x)\mathrm{\,d}x=\displaystyle\int  (3^x+2x)\mathrm{\,d}x=\dfrac{3^x}{\ln 3}+x^2+C.\]
	} 
\end{ex}

% nguồn: Phan1_Lop_12_Dai_so_Chuong_4
\begin{ex}%[2D4H1-4]%[ma:071a646b]
	Gọi $F(x)$ là một nguyên hàm của hàm số $f(x)=2x-\dfrac{1}{x}$ trên $(0;+\infty)$ thỏa mãn $F(1)=1$. 
	Tính $F(\mathrm{e})$. 
	\choice 
		{\True  $F(\mathrm{e})=\mathrm{e}^2-1$}
		{$F(\mathrm{e})=\mathrm{e}$} 
		 {$F(\mathrm{e})=\dfrac{1}{\mathrm{e}}$}
		{$F(\mathrm{e})=1-\mathrm{e}^2$} 
	\loigiai{Hàm số $F(x)$ là một nguyên hàm của $f(x)$ trên $(0;+\infty)$ thì $F(x)=x^2-\ln x  +C$, với $C$ là hằng số cần tìm.\\
		Theo đề  $F(1)=1\Leftrightarrow 1^2-\ln 1 +C =1\Leftrightarrow C=0$.\\
		Suy ra $F(x)=x^2-\ln x$.\\ 
	Vậy $F(\mathrm{e})=\mathrm{e}^2-\ln\mathrm{e} =\mathrm{e}^2-1$.  }
\end{ex}


%%%%%%%%%%%%%%%% ĐÚNG/SAI %%%%%%%%%%%%%%%%

% nguồn: Phan1_Lop_12_Dai_so_Chuong_4
\begin{ex}%[2D4H1-4]%[ma:b1f54580]
	Gọi $F(x)$ là một nguyên hàm của hàm $f(x)=\mathrm{e}^x-2x$ trên $\mathbb{R}$ thỏa mãn  $F(0)=1$.
	\choiceTF 
		{$F'(0)=0$} 
		{\True  $F(1)=\mathrm{e}-1$}
		{\True $\displaystyle\int F(x)\mathrm{\,d}x=\mathrm{e}^x-\dfrac{x^3}{3}+C$} 
		{$\displaystyle\int\dfrac{f(x)}{x\mathrm{e}^x}\mathrm{\,d}x=\ln|x|-2\mathrm{e}^x+C$} 
	 \loigiai{
		\begin{itemchoice}
			\itemch Ta có \[F'(x)=f(x)=\mathrm{e}^x-2x  \Rightarrow F'(0)=\mathrm{e}^0-0=1.\]
			\itemch Ta có \[F(x)=\displaystyle\int f(x)\mathrm{\,d}x=\mathrm{e}^x-x^2+C,\,F(0)=1 \Rightarrow  C=0.\]
			Do đó  $F(1)=\mathrm{e}-1$.
			 \itemch Ta có \[\displaystyle\int F(x)\mathrm{\,d}x=\int(\mathrm{e}^x-x^2)\mathrm{\,d}x=\mathrm{e}^x-\dfrac{x^3}{3}+C.\]
			\itemch Ta có  \[\displaystyle\int\dfrac{\mathrm{e}^x-2x}{x\mathrm{e}^x}\mathrm{\,d}x=\int\left(\dfrac{1}{x}-2\mathrm{e}^{-x}\right)\mathrm{\,d}x=\ln|x|+2\mathrm{e}^{-x}+C.\]
		\end{itemchoice}
	} 
\end{ex}

% nguồn: Phan1_Lop_12_Dai_so_Chuong_4
\begin{ex}%[2D4H1-4]%[ma:d7fd0baf]
	Biết $F(x)$ là một nguyên hàm của $f(x)=2x+\mathrm{e}^x$ thỏa mãn $F(0)=2\,025$.  Khi đó
	 \choiceTF
		 {$f(2)=4+\mathrm{e}$}
		{\True  $\displaystyle\int f(x)\mathrm{\,d}x=\displaystyle\int\left(2x+\mathrm{e}^x\right)\mathrm{\,d}x=x\mathrm{e}^2+\mathrm{e}^x+C$}
		{\True  $F(x)=x^2+\mathrm{e}^x+2\,024$}
		{$\displaystyle\int xf'(x^2)\mathrm{\,d}x=\displaystyle\int x\left(2+\mathrm{e}^{x^2}\right) \mathrm{\,d}x=x^2+x\mathrm{e}^{x^2}+C$} 
	 \loigiai{
		\begin{itemchoice}
			\itemch Ta có \[f(2)=2\cdot  2+\mathrm{e}^2=4+\mathrm{e}^2.\]
			\itemch Ta có \[\displaystyle\int  f(x)\mathrm{\,d}x=\displaystyle\int\left(2x+\mathrm{e}^x\right)\mathrm{\,d}x=x^2+\mathrm{e}^x+C.\]
			 \itemch Ta có $F(x)$ là một nguyên hàm của $f(x)$ và $F(0)=2\,025$.\\
			Suy ra \[\heva{&F(x)=x^{2}+\mathrm{e}^{x}+C\\&F(0)=2\,025} \Rightarrow 1+C=2\,025 \Leftrightarrow C=2\,024.\] 
			Vậy  $F(x)=x^2+\mathrm{e}^x+2\,024$.
			\itemch Ta có \[\displaystyle\int xf'\left(x^2\right) \mathrm{\,d}x=\displaystyle\int x\left(2+\mathrm{e}^{x^2}\right)\mathrm{\,d}x=x^2+\dfrac{1}{2} \mathrm{e}^{x^2}+C.\] 
		\end{itemchoice}
	} 
\end{ex}

% nguồn: Phan1_Lop_12_Dai_so_Chuong_4
\begin{ex}%[2D4H1-4]%[ma:7d8d10f6]
	Cho hàm số $f(x)=\mathrm{e}^x(2+\mathrm{e}^x)$ và $F(x)$ là nguyên hàm của  $f(x)$.
	 \choiceTF
		{\True $\displaystyle\int  f'(x)\mathrm{\,d}x=2\mathrm{e}^x+\mathrm{e}^{2x}+C$}
		{\True $F(x)=2\mathrm{e}^x+\dfrac{1}{2}\mathrm{e}^{2x}+C$} 
		{Nếu $F(0)=3$ thì  $F(1)=5\mathrm{e}$}
		{$F(0)=\dfrac{5}{2}$ thì phương trình $F(x)=0$ có  hai nghiệm phân biệt}
	 \loigiai{
		\begin{itemchoice}
			\itemch Ta có  \[\displaystyle\int f'(x)\mathrm{\,d}x=f(x)+C=2\mathrm{e}^x+\mathrm{e}^{2x}+C.\]
			\itemch Ta có \[F(x)=\displaystyle\int  f(x)\mathrm{\,d}x=\displaystyle\int(2\mathrm{e}^x+\mathrm{e}^{2x})\mathrm{\,d}x=2\mathrm{e}^x+\dfrac{1}{2}\mathrm{e}^{2x}+C.\]
			\itemch  Với $F(0)=3 \Rightarrow C=3-2-\dfrac{1}{2}=\dfrac{1}{2}$.\\
			Khi đó  \[F(x)=2\mathrm{e}^x+\dfrac{1}{2}\mathrm{e}^{2x}+\dfrac{1}{2}.\]
			Suy ra \[F(1)=2\mathrm{e}+\dfrac{1}{2}\mathrm{e}^2+\dfrac{1}{2}=\dfrac{\mathrm{e}^2+4\mathrm{e}+1}{2}.\] 
			\itemch Với $F(0)=\dfrac{5}{2} \Rightarrow C=0$,  suy ra $F(x)=2\mathrm{e}^x+\dfrac{1}{2}\mathrm{e}^{2x}$.\\
			Phương trình  $F(x)=0 \Leftrightarrow 2\mathrm{e}^x+\dfrac{1}{2}\mathrm{e}^{2x}=0$ vô nghiệm.
		\end{itemchoice}
	} 
\end{ex}

% nguồn: Phan1_Lop_12_Dai_so_Chuong_4
\begin{ex}%[2D4H1-4]%[ma:156a4671]
	Cho hàm số $f(x)=\mathrm{e}^x+x^\mathrm{e}$. 
	 \choiceTF
		{\True Nếu $F(x)$ là một nguyên hàm  của $f(x)$ thì $F(x)+5$ cũng là một nguyên hàm của $f(x)$}
		{$F(x)=\mathrm{e}^x+x^\mathrm{e}+C$ là họ nguyên  hàm của $f(x)$}
		{\True Nếu $F(x)$ là một  nguyên hàm của $f(x)$ và $F(0)=1$ thì $F(1)=2\mathrm{e}$}
		{Nếu $F(x)$, $G(x)$ lần  lượt là các nguyên hàm của $f(x)$ và $F(0)=1$, $G(1)=\mathrm{e}$ thì ta luôn có $G(x)=F(x)+\mathrm{e}$}
	\loigiai{ 
		\begin{itemchoice}
			\itemch Do  $\Big(F(x)+5\Big)'=F'(x)=f(x)$.
			\itemch  Ta có \[\displaystyle\int f(x)\mathrm{\,d}x=\displaystyle\int(\mathrm{e}^x+x^\mathrm{e})\mathrm{\,d}x=\mathrm{e}^x+\mathrm{e} x^{\mathrm{e}-1}+C.\]
			\itemch $F(x)=\mathrm{e}^x+\mathrm{e} x^{\mathrm{e}-1}$, với $F(0)=1 \Rightarrow F(1)=2\mathrm{e}$. 
			\itemch Từ điều kiện $F(0)=1$, $G(1)=\mathrm{e}$ suy ra  $G(x)=F(x)-\mathrm{e}$.
		\end{itemchoice}
	 }
\end{ex}

% nguồn: Phan1_Lop_12_Dai_so_Chuong_4
\begin{ex}%[2D4H1-4]%[ma:f62b56d9]
	Cho  $f(x)=\dfrac{1}{\sqrt[9]{x^7}};\,x\in (0;+\infty)$.
	\choiceTF 
		{$f(x)=\dfrac{1}{x^{\tfrac{9}{7}}};\,x\in (0;+\infty)$} 
		{\True  $f(x)={{x}^{-\tfrac{7}{9}}};\,x\in (0;+\infty)$}
		{\True $\forall \alpha \in (0;+\infty);\,\displaystyle\int{x^{\alpha }} \mathrm{\,d}x=\dfrac{x^{\alpha  +1}}{\alpha +1}+C$}
		 {$\displaystyle\int{f(x)}\mathrm{\,d}x=\dfrac{9}{2}\sqrt[9]{x^2}$}
	\loigiai{ 
		\begin{itemchoice}
			\itemch Sai. $f(x)=\dfrac{1}{x^{\tfrac{7}{9}}};\,x\in (0;+\infty)$. 
			\itemch Đúng. $f(x)=\dfrac{1}{x^{\tfrac{7}{9}}};\,x\in (0;+\infty)={{x}^{-\tfrac{7}{9}}};\,x\in (0;+\infty)$. 
			\itemch Đúng. 
			\itemch Sai. $\displaystyle\int{f(x)}\mathrm{\,d}x= \displaystyle\int{\dfrac{1}{\sqrt[9]{x^7}}}\mathrm{\,d}x= \displaystyle\int x^{-\tfrac{7}{9}}\mathrm{\,d}x =\dfrac{x^{-\tfrac{7}{9}+1}}{-\tfrac{7}{9}+1} + C=\dfrac{x^{\tfrac{2}{9}}}{\tfrac{2}{9}}  + C=\dfrac{9}{2}\sqrt[9]{x^2} + C$.
		\end{itemchoice}
	 }
\end{ex}

% nguồn: Phan1_Lop_12_Dai_so_Chuong_4
\begin{ex}%[2D4H1-4]%[ma:7dc95952]
	Cho hàm số $F(x)$ là một nguyên hàm  của hàm số $f(x)=\dfrac{6^{x+1}+3^{2x}}{2^{2x+3}}$.
	 \choiceTF
		 {\True $F'(x)=\dfrac{6^{x+1}+3^{2x}}{2^{2x+3}}$}
		 {$F(x)=\dfrac{\displaystyle\int(6^{x+1}+3^{2x})\mathrm{\,d}x}{\displaystyle\int 2^{2x+3}\mathrm{\,d}x}$}
		{\True $F(x)=\displaystyle\int\left[\dfrac{3}{4}\left(\dfrac{3}{2}\right)^x+\dfrac{1}{8}\left(\dfrac{9}{4}\right)^x\right]\mathrm{\,d}x$} 
		{\True Nếu $F(0)=0$  thì $F(1)=\dfrac{29}{64(\ln 3-\ln 2)}$}
	\loigiai{ 
		\begin{itemchoice}
			\itemch Theo  định nghĩa nguyên hàm suy ra $F'(x)=f(x)=\dfrac{6^{x+1}+3^{2x}}{2^{2x+3}}$.
			\itemch Theo tính chất  nguyên hàm suy ra $F(x)=\displaystyle\int\dfrac{6^{x+1}+3^{2x}}{2^{2x+3}}\mathrm{\,d}x \ne \dfrac{\displaystyle\int(6^{x+1}+3^{2x})\mathrm{\,d}x}{\displaystyle\int 2^{2x+3}\mathrm{\,d}x}$.
			\itemch Ta có \[F(x)=\displaystyle\int\dfrac{6^{x+1}+3^{2x}}{2^{2x+3}}\mathrm{\,d}x=\int\dfrac{6\cdot 6^x+9^x}{8\cdot 4^x}\mathrm{\,d}x=\displaystyle\int\left[\dfrac{3}{4}\left(\dfrac{3}{2}\right)^x+\dfrac{1}{8}\left(\dfrac{9}{4}\right)^x\right]\mathrm{\,d}x.\] 
			\itemch Ta có  \[F(x)=\dfrac{\dfrac{3}{4}\left(\dfrac{3}{2}\right)^x+\dfrac{1}{16}\left(\dfrac{3}{2}\right)^{2x}}{\ln 3-\ln 2}+C.\]
			Với $F(0)=0 \Rightarrow  C=-\dfrac{13}{16(\ln 3-\ln 2)}$.\\
			Suy ra 
			\[F(1)=\dfrac{81}{64(\ln 3-\ln 2)}-\dfrac{52}{64(\ln 3-\ln 2)}=\dfrac{29}{64(\ln 3-\ln  2)}.\]
		\end{itemchoice}
	} 
\end{ex}

% nguồn: Phan1_Lop_12_Dai_so_Chuong_4
\begin{ex}%[2D4H1-4]%[ma:99bf3977]
	Cho hàm số  $f(x)$ xác định trên khoảng $K$. Gọi $F(x)$ là họ nguyên hàm của $f(x)$ trên $K$.
	\choiceTF 
		{\True $F'(x) =  f(x)$}
		 {Nếu $f(x) = \dfrac{1}{x}$ thì $F(x) = \ln x + C$}
		{Nếu $f(x) = \cos x$ thì $F(x)  = \sin x$}
		{\True Nếu $f(x) = \mathrm{e}^x$ thì $F(x) = \mathrm{e}^x +  C$}
	\loigiai{\begin{itemchoice}
	\itemch  \textbf{Đúng}. Định nghĩa nguyên hàm $F'(x) = f(x)$.
	 \itemch \textbf{Sai}. Nguyên hàm của $\dfrac{1}{x}$ là $\ln|x| + C$ chứ không phải $\ln x + C$ với mọi $x$.
	\itemch \textbf{Sai}. $\displaystyle\int \cos x \,\mathrm{\,d}x = \sin x + C$ nên thiếu hằng số  $C$.
	\itemch \textbf{Đúng}.  $\displaystyle\int \mathrm{e}^x \mathrm{d}x = \mathrm{e}^x + C$.
\end{itemchoice}
} 
\end{ex}

% nguồn: Phan1_Lop_12_Dai_so_Chuong_4
\begin{ex}%[2D4H1-4]%[ma:84aa2e85]
	Cho  hàm số $F(x)=\displaystyle\int\left(4^x+2^x+1\right)\mathrm{\,d}x$.
	\choiceTF 
		{\True  $F(x)=\displaystyle\int 4^x\mathrm{\,d}x+\displaystyle\int 2^x\mathrm{\,d}x+\displaystyle\int\mathrm{\,d}x$}
		{$F(x)=\dfrac{2^{2x}+2^x}{\ln  2}+x+C$}
		{Nếu $F(1)=\dfrac{4}{\ln 2}$ thì  $F(2)=\dfrac{12}{\ln 2}-1$}
		{Nếu $G(x)=\displaystyle\int \dfrac{8^x+1}{2^x+1}\mathrm{\,d}x$ thì  $G(x)=F(x)+C$}
	 \loigiai{
		\begin{itemchoice}
			\itemch Ta có \[F(x)=\displaystyle\int\left(4^x+2^x+1\right)\mathrm{\,d}x=\displaystyle\int 4^x\mathrm{\,d}x+\displaystyle\int 2^x\mathrm{\,d}x+\displaystyle\int\mathrm{\,d}x.\] 
			\itemch Ta có \[F(x)=\displaystyle\int\left(4^x+2^x+1\right)\mathrm{\,d}x=\dfrac{4^x}{\ln 4}+\dfrac{2^x}{\ln 2}+x+C=\dfrac{2^{2x}}{2\ln 2}+\dfrac{2^x}{\ln  2}+x+C=\dfrac{2^{2x-1}+2^x}{\ln 2}+x+C.\]
			\itemch Thay $x=1$  vào $F(x)$ ta được
			 \[F(1)=\dfrac{2^{1}+2^1}{\ln 2}+1+C=\dfrac{4}{\ln 2}+1+C=\dfrac{4}{\ln 2} \Rightarrow C=-1.\]
			 Suy ra $F(x)=\dfrac{2^{2x-1}+2^x}{\ln 2}+x-1$.\\
			 Suy ra $F(2)=\dfrac{2^{3}+2^2}{\ln 2}+2-1=\dfrac{12}{\ln 2}+1$.
			\itemch Ta có \[G(x)=\dfrac{8^x+1}{2^x+1}=\dfrac{2^{3x}+1}{2^x+1}=\dfrac{(2^x+1)(4^x-2^x+1)}{2^x+1}=4^x-2^x+1.\] 
		\end{itemchoice}
	} 
\end{ex}

% nguồn: Phan1_Lop_12_Dai_so_Chuong_4
\begin{ex}%[2D4H1-4]%[ma:f5597647]
	Cho hàm số $f(x)=\dfrac{2^{2x+1}-1+\cos 2x}{2^x+\sin x}$ trên $[0;+\infty)$  thỏa mãn $F(0)=\dfrac{2}{\ln 2}+2$.
	 \choiceTF
		{$f(\pi)=2^{\pi+1}-1$} 
		 {\True $F'(0)=0$}
		{\True Hệ số góc của tiếp tuyến với đồ thị hàm số $F(x)$ tại điểm có  hoành độ $x_0=\dfrac{\pi}{2}$ là $k=2\sqrt{2^\pi}-2$}
		{$F(-1)=\dfrac{1}{\ln 2}+2\cos  1$}
	 \loigiai{
		\begin{itemchoice}
			\itemch Ta có  \[f(x)=\dfrac{2^{2x+1}-1+\cos 2x}{2^x+\sin x}=\dfrac{2(2^x-\sin x)(2^x+\sin x)}{2^x+\sin x}=2\cdot 2^x-2\sin x.\]
			 Suy ra $f(\pi)=2\cdot 2^{\pi}-2\sin \pi=2^{\pi+1}$.
			\itemch Ta có \[F'(x)=f(x)=2\cdot 2^x-2\sin  x \Rightarrow F'(0)=2\cdot 2^0-2\cdot 0=2.\]
			\itemch Hệ số góc của tiếp tuyến với đồ thị hàm số $F(x)$ tại  điểm có hoành độ $x_0=\dfrac{\pi}{2}$ là \[k=F'\left(\dfrac{\pi}{2}\right)=f\left(\dfrac{\pi}{2} \right)={2\cdot 2}^{\tfrac{\pi}{2}}-2\sin \dfrac{\pi}{2}=2\cdot \sqrt{2^{\pi}}-2.\]
			\itemch Nguyên hàm  của $f(x)$ là $F(x)=\dfrac{2^{x+1}}{\ln 2}+2\cos x+C$, với $F(0)=\dfrac{2}{\ln 2}+2 \Rightarrow C=0$.\\
			Suy  ra $F(-1)=\dfrac{1}{\ln 2}+2\cos 1$.
		\end{itemchoice}
	 }
\end{ex}

% nguồn: Phan1_Lop_12_Dai_so_Chuong_4
\begin{ex}%[2D4H1-4]%[ma:f8405c08]
	Cho hàm số  $f(x)=2^{-3x+1}$.
	 \choiceTF
		{\True $f(x)=2\left(\dfrac{1}{8}  \right)^x$}
		{$\displaystyle\int{f(x)\mathrm{\,d}x}=\dfrac{1}{2}\displaystyle\int{\left(\dfrac{1}{8} \right)^x  \mathrm{\,d}x}$}
		 {$\displaystyle\int{{{\left(\dfrac{1}{8} \right)}^x}\mathrm{\,d}x}=\dfrac{\left(\dfrac{1}{8} \right)^x}{\ln\left(\dfrac{1}{8} \right)}$}
		{\True $\displaystyle\int {f(x)\mathrm{\,d}x}=\dfrac{2^{-3x+1}}{-3\ln 2}+C$} 
	 \loigiai{
		\begin{itemchoice}
			\itemch Đúng. $f(x)=2^{-3x+1}  = 2\cdot 2^{-3x} = 2\left(\dfrac{1}{8}\right)^{x}$.
			\itemch Sai. $\displaystyle\int{f(x)\mathrm{\,d}x} = \int{2\left(\dfrac{1}{8}\right)^{x}\mathrm{\,d}x}  = 2\displaystyle\int{\left(\dfrac{1}{8} \right)^x \mathrm{\,d}x}$.
			\itemch Sai. $\displaystyle\int{{{\left(\dfrac{1}{8} \right)}^x}\mathrm{\,d}x}=\dfrac{\left(\dfrac{1}{8}  \right)^x}{\ln\left(\dfrac{1}{8} \right)} + C$.
			\itemch  Đúng. $\displaystyle\int{f(x)\mathrm{\,d}x} = \int{2\left(\dfrac{1}{8}\right)^{x}\mathrm{\,d}x} = 2\displaystyle\int{\left(\dfrac{1}{8} \right)^x \mathrm{\,d}x} = 2 \dfrac{\left(\dfrac{1}{8}\right)^x}{\ln\left(\dfrac{1}{8}\right)} + C = \dfrac{2^{-3x+1}}{-3\ln 2} + C$.
		\end{itemchoice}
	 }
\end{ex}

% nguồn: Phan1_Lop_12_Dai_so_Chuong_4
\begin{ex}%[2D4V1-4]%[ma:3dd7c7cf]
	Cho hàm số $f(x)=3^x$. Khi đó 
	 \choiceTF
		{$\displaystyle\int  f(x)\mathrm{\,d}x=3^x \ln 3+C$}
		{\True $\displaystyle\int  \left[f(x)-\mathrm{e}^x \right]\mathrm{\,d}x=\dfrac{3^x}{\ln 3}-\mathrm{e}^x+C$}
		{\True $\displaystyle\int f(x).\mathrm{e}^x \mathrm{\,d}x=\dfrac{3^x \cdot  \mathrm{e}^x}{1+\ln 3}+C$}
		{\True Nếu $F(x)$ là một nguyên hàm của $f(x)$ thì $F'\left(\log_9 5\right)=\sqrt{5}$} 
	\loigiai{ 
		\begin{itemchoice}
			\itemch {\bf Sai}.  \\
			$\displaystyle\int f(x)\mathrm{\,d}x=\displaystyle\int 3^x \mathrm{\,d}x=\dfrac{3^x}{\ln x}+C$. 
			\itemch {\bf Đúng}.  \\
			$\displaystyle\int  \left[f(x)-\mathrm{e}^x \right]\mathrm{\,d}x=\displaystyle\int f(x)\mathrm{\,d}x-\displaystyle\int \mathrm{e}^x \mathrm{\,d}x=\dfrac{3^x}{\ln 3}-\mathrm{e}^x+C$.
			 \itemch {\bf Đúng}.\\
			 $\displaystyle\int f(x)\cdot \mathrm{e}^x \mathrm{\,d}x=\displaystyle\int (3\mathrm{e})^x \mathrm{\,d}x=\dfrac{(3 \mathrm{e})^x}{\ln(3 \mathrm{e})}+C=\dfrac{3^x \cdot \mathrm{e}^x}{1+\ln 3}+C$.
			\itemch {\bf Đúng}. \\ 
			Ta có $F'(x)=f(x)$ suy  ra $F'\left(\log_9 5\right)=f\left(\log_9 5\right)=3^{\log_9 5}=3^{\frac{1}{2} \log_3 5}=\sqrt{5}$.
		\end{itemchoice}
	} 
\end{ex}


%%%%%%%%%%%%%%%% TRẢ LỜI NGẮN %%%%%%%%%%%%%%%%

% nguồn: Phan2_Lop_12_Dai_so_Chuong_4
\begin{bt}%[2D4H1-4]%[ma:13fda258]
	Giả sử $\displaystyle\int \dfrac{1}{\sqrt[3]{x^5}}\, \mathrm{\,d}x=ax^{\tfrac{m}{n}}+C$ với $a$  là hằng số thực, $n$ là số nguyên dương, $m$ là số nguyên và ước số chung lớn nhất của $m$ và $n$ bằng $1$. Giá trị của biểu thức $S=a+m+n$ là bao nhiêu?
	 \shortans{-0,5}
	 \loigiai{
		Ta có $\begin{aligned}
		\displaystyle\int \dfrac{1}{\sqrt[3]{x^5}}\mathrm{\,d}x&=\displaystyle\int  \dfrac{1}{x^{\tfrac{5}{3}}}\mathrm{\,d}x=\displaystyle\int x^{\tfrac{-5}{3}}\mathrm{\,d}x=\dfrac{x^{\tfrac{-5}{3}+1}}{\dfrac{-5}{3}+1}+C\\
		&=\dfrac{-3}{2} \dfrac{1}{x^{\tfrac{2}{3}}}+C=\dfrac{-3}{2} x^{\tfrac{-2}{3}}+C. 
	\end{aligned}$\\
	Suy ra $a+m+n=\dfrac{-3}{2}-2+3=-0{,}5$. 
} 
\end{bt}

% nguồn: Phan2_Lop_12_Dai_so_Chuong_4
\begin{bt}%[2D4H1-4]%[ma:fb372049]
	Cho hàm số $f(x)=-(2x^2-8x+7)\mathrm{e}^{-2x}$ có một nguyên hàm là $F(x)=(ax^2+bx+c)\mathrm{e}^{-2x}$.  Tính $a+b+c$.
	\shortans{0} 
	 \loigiai{
		Ta có 
		\[F'(x)=(2ax+b)\mathrm{e}^{-2x}-2(ax^2+bx+c)\mathrm{e}^{-2x}=-\left[2ax^2-(2a-2b)x-b+2c\right]\mathrm{e}^{-2x}.\] 
		Do $F'(x)=f(x)$  nên
		\[\heva{&2a=2\\&-2a+2b=-8\\&-b+2c=7} \Leftrightarrow  \heva{&a=1\\&b=-3\\&c=2.}\]
		 Vậy $a+b+c=0$.
	} 
\end{bt}

% nguồn: Phan2_Lop_12_Dai_so_Chuong_4
\begin{bt}%[2D4H1-4]%[ma:19dba95d]
	Cho $F(x)$ là một nguyên hàm của hàm số $f(x)=\mathrm{e}^{x}+2x$ thỏa mãn $F(0)=\dfrac{3}{2}$. Tìm $F(2)$ (kết  quả làm tròn đến hàng phần chục).
	 \shortans{11,9}
	\loigiai{ 
		Ta có  \[F(x)=\displaystyle\int (\mathrm{e}^x+2x)\mathrm{\,d}x=\mathrm{e}^x+x^2+C.\]
		Từ $F(0)=\dfrac{3}{2}$ suy ra $\mathrm{e}^0+C=\dfrac{3}{2} \Leftrightarrow  C=\dfrac{1}{2}$.\\
		 Suy ra $F(x)=\mathrm{e}^x+x^2+\dfrac{1}{2}$.\\
		Vậy  $F(2)=\mathrm{e}^2+4+\dfrac{1}{2}\approx 11{,}9$.
	 }
\end{bt}

% nguồn: Phan2_Lop_12_Dai_so_Chuong_4
\begin{bt}%[2D4H1-4]%[ma:aafcda5f]
	Biết $F(x)=\mathrm{e}^{2x}+2x^2$ là một nguyên  hàm của hàm số $f(x)$ trên $\mathbb{R}$. Giả sử $G(x)$ là một nguyên hàm của hàm số $f(2x)+\cos \pi x$. Khi đó $G(2)-G(0)=a\cdot\mathrm{e}^b+c$ với $a$, $b$, $c\in\mathbb{Q}$. Tính $abc$.
	 \shortans{62}
	 \loigiai{
		Ta có \[f(x)=F'(x)=2\mathrm{e}^{2x}+4x \Rightarrow f(2x)=2\mathrm{e}^{4x}+8x.\] 
		Do đó 
		\[G(x)=\displaystyle\int [f(2x)+\cos \pi x]\mathrm{\,d}x=\displaystyle\int \left(2\mathrm{e}^{4x}+8x+\cos \pi  x\right)\mathrm{\,d}x=\dfrac{1}{2}\mathrm{e}^{4x}+4x^2+\dfrac{1}{\pi}\sin \pi x+C.\]
		Suy ra  \[G(2)-G(0)=\left(\dfrac{1}{2}\mathrm{e}^8+16+C\right)-\left(\dfrac{1}{2}+C\right)=\dfrac{1}{2}\mathrm{e}^8+\dfrac{31}{2}.\]
		Vậy $abc=62$. 
	} 
\end{bt}

% nguồn: Phan2_Lop_12_Dai_so_Chuong_4
\begin{bt}%[2D4H1-4]%[ma:5bae8a84]
	Cho $F(x)=\left(ax^2+bx+c\right)\mathrm{e}^{2x}$ là một nguyên hàm của hàm số $f(x)=2x^2\cdot  \mathrm{e}^{2x}$. Tính tổng $a+b+2c$.
	\shortans{1} 
	\loigiai{ 
		Ta  có \[F'(x)=(2ax+b)\mathrm{e}^{2x}+2(ax^2+bx+c)\mathrm{e}^{2x}=\left[2ax^2+(2a+2b)x+b+2c\right]\mathrm{e}^{2x}.\]
		Mà $F'(x)=f(x)=2x^2\mathrm{e}^{2x}$, $\forall  x\in \mathbb{R}$ nên suy ra
		\[\heva{&2a=2\\&2a+2b=0\\&b+2c=0}\Leftrightarrow \heva{&a=1\\&b=-1\\&c=\dfrac{1}{2}.}\] 
		Vậy $a+b+2c=1$. 
	} 
\end{bt}


%%%%%%%%%%%%%%%% TỰ LUẬN %%%%%%%%%%%%%%%%

% nguồn: Phan2_Lop_12_Dai_so_Chuong_4
\begin{bt}%[2D4H1-4]%[ma:bfa1e0ed]
	Tính đạo hàm của $F(x)=\ln(x+\sqrt{x^2+1})$. Từ đó suy  ra nguyên hàm của $f(x)=\dfrac{1}{\sqrt{x^2+1}}$.
	 \loigiai{
		Ta có $f(x)=F'(x)=\dfrac{1+\dfrac{x}{\sqrt{x^2+1}}}{x+\sqrt{x^2+1}}=\dfrac{\dfrac{x+\sqrt{x^2+1}}{\sqrt{x^2+1}}}{x+\sqrt{x^2+1}}=\dfrac{1}{\sqrt{x^2+1}}$.\\ 
		Vậy $\displaystyle\int \dfrac{1}{\sqrt{x^2+1}}\mathrm{\,d}x=\ln(x+\sqrt{x^2+1})$. 
	} 
\end{bt}

% nguồn: Phan2_Lop_12_Dai_so_Chuong_4
\begin{bt}%[2D4H1-4]%[ma:e5c11ac6]
	Cho hàm số $f(x)=\mathrm{e}^{2x}+x^2$. Tìm họ các nguyên hàm của hàm số  $g(x)=2f(x)$.
	\loigiai{ 
		Ta có $\displaystyle\int g(x) \mathrm{\,d}x=\displaystyle\int 2f(x)  \mathrm{\,d}x=\displaystyle\int 2(\mathrm{e}^{2x}+x^2)\mathrm{\,d}x =2\displaystyle\int \mathrm{e}^{2x} \mathrm{\,d}x+2\displaystyle\int x^2 \mathrm{\,d}x=\mathrm{e}^{2x}+\dfrac{2}{3}x^3+C$.
	} 
\end{bt}

% nguồn: Phan2_Lop_12_Dai_so_Chuong_4
\begin{bt}%[2D4H1-4]%[ma:6969638d]
	Tìm  họ nguyên hàm của hàm số $f(x)$, biết:
	\begin{multicols}{3}
		\begin{enumerate}
			\item  $f(x)=2\,024^x$.
			\item  $f(x)=5^x+1$.
			 \item $f(x)=\mathrm{e}^{-x}$.
			\item  $f(x)=\mathrm{e}^{2x}-\dfrac{1}{x^2}$.
			\item $f(x)=8^x \cdot  2^{1-2x}$.
			 \item $f(x)=\left(3^x-\dfrac{1}{3^x}\right)^2$.
		\end{enumerate}
	\end{multicols}
	\loigiai{ 
		\begin{enumerate}
			\item $\displaystyle\int 2\,024^x \mathrm{\,d}x=\dfrac{2\,024^x}{\ln 2\,024}+C$. 
			\item $\displaystyle\int (5^x+1) \mathrm{\,d}x=\dfrac{5^x}{\ln  5}+x+C$.
			\item $\displaystyle\int \mathrm{e}^{-x}  \mathrm{\,d}x=-\mathrm{e}^{-x}+C$.
			\item $\displaystyle\int \left(\mathrm{e}^{2x}-\dfrac{1}{x^2}\right) \mathrm{\,d}x=\dfrac{1}{2}\mathrm{e}^{2x}+\dfrac{1}{x}+C$. 
			\item Ta có $f(x)=8^x \cdot 2^{1-2x}=2^{3x}  \cdot 2^{1-2x}=2^{x+1}$.\\
			Suy ra $\displaystyle \int f(x) \mathrm{\,d}x=\displaystyle\int  2^{x+1}\mathrm{\,d}x=\dfrac{2^{x+1}}{\ln 2}+C$.
			\item  Ta có $f(x)=\left(3^x-\dfrac{1}{3^x}\right)^2=3^{2x}-2+3^{-2x}=9^x-2+9^{-x}$.\\
			Suy ra $\displaystyle \int f(x) \mathrm{\,d}x=\dfrac{9^x}{\ln  9}-2x-\dfrac{1}{9^x\ln 9}+C$.
		\end{enumerate}
	} 
\end{bt}

% nguồn: Phan2_Lop_12_Dai_so_Chuong_4
\begin{bt}%[2D4H1-4]%[ma:9fa42e93]
	Tìm họ nguyên hàm của các hàm số sau: 
	\begin{multicols}{3}
		\begin{enumerate}
			\item $y=2^x$. 
			 \item $y=\dfrac{1}{3^x}$.
			 \item $y=2\mathrm{e}^x-5^x$.
			\item  $y=\dfrac{1}{\mathrm{e}^x}$.
			 \item $y=2\cdot 3^x-\dfrac{1}{3}\cdot 7^x$.
			\item  $y=2^{3x} \cdot 3^{2x}$.
			 \item $y=\dfrac{2^{x+1}-5^{x-1}}{10^x}$.
			\item $y=\dfrac{\mathrm{e}^{3x}+1}{\mathrm{e}^x+1}$. 
			\item  $y=\mathrm{e}^x\left(2017-\dfrac{2018\mathrm{e}^{-x}}{x^5}\right)$.
		\end{enumerate}
	\end{multicols}
	\loigiai{ 
		\begin{enumerate}
			\item  Ta có \[\displaystyle \int 2^x \mathrm{\,d}x=\dfrac{2^x}{\ln 2}+C.\]
			\item  Ta có \[\displaystyle \int \left(\dfrac{1}{3}\right)^x \mathrm{\,d}x=-\dfrac{1}{3^x\ln 3}+C.\]
			\item Ta có \[\displaystyle \int (2\mathrm{e}^x-5^x) \mathrm{\,d}x=2\mathrm{e}^x-\dfrac{5^x}{\ln 5}+C.\] 
			 \item Ta có \[\displaystyle \int \dfrac{1}{\mathrm{e}^x} \mathrm{\,d}x=\int \mathrm{e}^{-x}\mathrm{\,d}x=-\mathrm{e}^{-x}+C.\]
			 \item Ta có \[\displaystyle \int \left(2\cdot 3^x-\dfrac{1}{3}\cdot 7^x\right) \mathrm{\,d}x=2\dfrac{3^x}{\ln 3}-\dfrac{7^x}{3\ln 7}+C.\]
			\item Ta có \[y=8^x \cdot 9^x=72^x \Rightarrow \displaystyle \int y\mathrm{\,d}x=\dfrac{72^x}{\ln  72}+C.\]
			\item Ta có \[y=\dfrac{2^{x+1}-5^{x-1}}{10^x}=2\cdot \left(\dfrac{1}{5}\right)^x-\dfrac{1}{5}\cdot \left(\dfrac{1}{2}\right)^x.\] 
			 Suy ra
			\[\displaystyle  \int y \mathrm{\,d}x=-\dfrac{2}{5^x\ln 5}+\dfrac{1}{5\cdot 2^x\ln 2}+C.\]
			\item Ta có \[\dfrac{\mathrm{e}^{3x}+1}{\mathrm{e}^x+1}=\mathrm{e}^{2x}-\mathrm{e}^x+1 \Rightarrow \displaystyle \int y  \mathrm{\,d}x=\dfrac{1}{2}\mathrm{e}^{2x}-\mathrm{e}^x+x+C.\]
			 \item Ta có \[y=\mathrm{e}^x\left(2\,017-\dfrac{2\,018\mathrm{e}^{-x}}{x^5}\right)=2\,017\mathrm{e}^x-\dfrac{2\,018}{x^5}.\]
			Suy  ra
			\[\displaystyle \int y  \mathrm{\,d}x=2\,017\mathrm{e}^x+\dfrac{504\cdot 5}{x^4}+C.\]
		\end{enumerate}
	} 
\end{bt}

% nguồn: Phan2_Lop_12_Dai_so_Chuong_4
\begin{bt}%[2D4H1-4]%[ma:cf9643c9]
	Cho $f(x)=x^2\ln x$ và $g(x)=x\ln x$. Tính $f'(x)$ và  $\displaystyle\int g(x)\mathrm{\,d}x$.
	\loigiai{ 
		Ta có: $f'(x)=2x\ln x+\dfrac{1}{x}\cdot  x^2=2x\ln x+x$.\\
		Tính $I=\displaystyle\int g(x)  \mathrm{\,d}x=\displaystyle\int x\ln x\mathrm{\,d}x$.\\
		Đặt $\heva {&u=\ln  x\Rightarrow \mathrm{\,d}u=\dfrac{1}{x}\mathrm{\,d}x\\&\mathrm{\,d}v=x\Rightarrow v=\dfrac{x^2}{2}}$\\
		$\Rightarrow I=\dfrac{x^2}{2}\ln x-\displaystyle\int \dfrac{1}{x}\cdot\dfrac{x^2}{2}\mathrm{\,d}x=\dfrac{x^2}{2}\ln x-\displaystyle\int \dfrac{x}{2}\mathrm{\,d}x=\dfrac{x^2}{2}\ln x-\dfrac{x^2}{4}+C$. 
	 }
\end{bt}

% nguồn: Phan2_Lop_12_Dai_so_Chuong_4
\begin{bt}%[2D4H1-4]%[ma:6dae0660]
	Cho $F(x)$ là một nguyên hàm của hàm số $f(x)=3^x$ trên $\mathbb{R}$ thỏa mãn $F(0)=0$. Tính  $F(1)$.
	\loigiai{Ta có $$F(1)-F(0)=\displaystyle\int_0^1 f(x) \mathrm{d} x=\displaystyle\int_0^1 3^x \mathrm{\,d} x=\dfrac{3^x}{\ln 3}\bigg|_0  ^1=\dfrac{3}{\ln 3}-\dfrac{1}{\ln 3}=\dfrac{2}{\ln 3}.$$
		 Suy ra $F(1)=\dfrac{2}{\ln 3}+F(0)=\dfrac{2}{\ln 3}+0=\dfrac{2}{\ln 3}$.
	 }
\end{bt}
